% Source: Wald GR, physical PDF p. 26 (printed p. 13).
% Coverage: Figure 2.1, transition map for overlapping coordinate charts.
% Figures/tables/footnotes: Figure only; no tables or footnotes.
% Status: source-reviewed and compile-clean.
% Uncertainties: none.
\begingroup
\renewcommand{\thefigure}{2.1}
\begin{figure}[htbp]
  \centering
\begin{tikzpicture}[x=0.72cm,y=0.72cm,>=Stealth,font=\small]
  \draw[rounded corners=5pt] (0,3.4)
    .. controls (-0.3,4.4) and (0.4,5.1) .. (1.3,4.9)
    .. controls (2.0,4.1) and (2.5,4.8) .. (3.5,4.7)
    .. controls (4.7,4.6) and (5.1,3.8) .. (4.7,3.0)
    .. controls (5.0,1.7) and (3.7,1.5) .. (2.8,1.9)
    .. controls (1.8,1.5) and (0.2,2.0) .. (0,3.4);
  \node at (0.8,4.2) {\(M\)};
  \draw (1.25,3.2) ellipse (0.92 and 0.72);
  \draw (2.1,3.25) ellipse (0.64 and 0.78);
  \fill[gray!25] (1.72,3.18) ellipse (0.18 and 0.54);
  \node at (0.88,2.78) {\(\mathcal O_\alpha\)};
  \node at (2.32,2.70) {\(\mathcal O_\beta\)};

  \draw[->] (2.70,3.25) -- (7.62,3.25);
  \node[above=1pt] at (5.05,3.28) {\(\psi_\beta\)};
  \draw (6.55,2.0) rectangle (9.75,5.05);
  \node at (7.1,4.62) {\(\mathbb{R}^n\)};
  \draw (8.05,3.25) ellipse (0.75 and 0.66);
  \fill[gray!25] (7.65,3.25) ellipse (0.19 and 0.46);
  \node at (8.35,3.25) {\(U_\beta\)};

  \draw (0.30,-2.30) rectangle (3.70,0.75);
  \node at (0.82,0.32) {\(\mathbb{R}^n\)};
  \draw (2.05,-0.78) ellipse (0.78 and 0.67);
  \fill[gray!25] (2.68,-0.78) ellipse (0.19 and 0.47);
  \node at (1.82,-0.78) {\(U_\alpha\)};
  \draw[->] (1.32,2.48) -- (1.98,-0.24);
  \node[right=1pt] at (1.94,1.05) {\(\psi_\alpha\)};
  \draw[->] (2.80,-1.00) -- (7.70,2.88)
    node[pos=0.55,below,sloped] {\(\psi_\beta\circ\psi_\alpha^{-1}\)};
  % Reserve trim clearance above the artwork when this float lands at page top.
  \path[draw=none] (0,6.1) circle[radius=.1pt];
\end{tikzpicture}
  \caption{An illustration of the map \(\psi_\beta\circ\psi_\alpha^{-1}\)
    arising when two coordinate systems overlap.}
  \label{fig:2.1}
\end{figure}
\endgroup
